\begin{pastexam}[label={px:ncku-111-ee-ds-17}]{成大電機所 111 學年度 單選題第 17 題}
（單選；每題 4 分）

To represent the hierarchical relationship between elements, which data structure is suitable?

A.~Dequeue \qquad B.~Priority \qquad C.~Tree \qquad D.All of above
\end{pastexam}

\begin{pastsoln}
答案是 C。

階層關係（例如公司組織圖、檔案目錄：上層包含下層）用\textbf{樹}表示：父節點在上、子節點在下，每個節點只有一個父節點
（資結書第~7.1~節）。
\begin{itemize}
\item A.~雙端佇列：線性結構，元素排成一列，沒有上下層。 錯。
\item B.~優先（佇列）：只管「誰的優先權最高」，不表示元素之間的從屬關係。 錯。 （優先佇列常用堆積實作，堆積的形狀是樹，
但那是實作方式，不是用來表示階層關係。）
\item C.~樹：對。
\item D.~全部：A、B 不是。 錯。
\end{itemize}
\end{pastsoln}
